% TT14200C
% Fraktur
% 14.0
% 2013-11-30

.. _m-fraktur:


Taktik \Ca{Zeitkritisch. Evtl. vitale
Bedrohung v. a. durch Blutverlust.}
Erkennen
von großen Blutverlusten, Blutstillung,  Steriles
Abdecken, Einschätzen der vitalen Bedrohung, ggfs
Schmerztherapie. 
Bei *vitaler Bedrohung*: Zeitkritisch,
zügiger Transport. 
Bei *stabilen Patienten*: Schienung}




\begin{itemize}
  -    **Blutstillung**: Z. B. mit Beckengurt
  (Beckenfraktur)
  -   Vitale Bedrohung einschätzen, ggfs. |TxMassVitMK|.
  Ggfs. mit hypovolämischen Schock rechnen, Schockbekämpfung
  -    **Wundversorgung**: Bei offenen Frakturen ist
  Keimfreiheit be\-sonders wichtig, da  die Gefahr einer lebenslangen
  Knochen\-marks\-ent\-zündung \Z{(Osteomyelitis)} besteht!
  Wenn möglich soll eine *Wundfolie* (z. B.
  \index{OpSite}\T{OpSite}\texttrademark ) verwendet werden.
  -   Kleidung und Schmuck von den betroffenen Gliedmaßen
  entfernen (Schwellung!)
  -   Ggfs. Schmerztherapie (veranlassen)
  -   Ruhigstellung/Schienung bei stabilen Patienten
  -   Beurteilung von Durchblutung, Motorik und
  Sensibilität (DMS) an der betroffenen Extremität und an
  der Gegenseite
\end{itemize}